%! AP Statistics — Unit 7, Lesson 7.2 — Group Activity
\documentclass[10pt]{article}
\usepackage{apstats-article}
\usepackage{apstats-boxes}
\newcommand{\TallMath}[1]{$\displaystyle #1\rule[-1.4em]{0pt}{3.2em}$}

\begin{document}

\pageheader{Unit 7, Lesson 7.2}{Group Activity}
\namepartnerperiod

\vspace{0.1in}
\textbf{Directions:} For each problem, use the full four-step procedure (Parameter,
Conditions, Calculate, Conclude). Show all work.

\begin{tcolorbox}[colback=white, colframe=black!40, title=\textbf{Tier R --- Remediate},
    fonttitle=\bfseries, arc=2mm, left=3mm, right=3mm, top=2mm, bottom=2mm]

\textbf{Sleep Study.} A school counselor wants to estimate the mean nightly sleep
duration for students at her school. A random sample of 25 students yields
$\bar x = 6.8$ hours and $s = 1.2$ hours.

Construct a 90\% confidence interval for the true mean nightly sleep duration.

\textbf{Step 1 --- Parameter:} $\mu = $ \writeline

\textbf{Step 2 --- Conditions:} (conditions are pre-verified for you)
\begin{itemize}
    \item Random: random sample of students. \checkmark
    \item 10\%: $25 \le 10\%$ of all students at the school. \checkmark
    \item Large Sample: $n = 25 < 30$; assume counselor verified no strong skewness or
          outliers in the data. \checkmark
\end{itemize}

\textbf{Step 3 --- Calculate:}

$df = $ \blank{2cm} \quad $t^* = $ \blank{2.5cm} (from $t$-table, 90\% CI)

$SE = s/\sqrt{n} = 1.2/\sqrt{25} = $ \blank{2.5cm}

$ME = t^* \cdot SE = $ \blank{4cm}

Interval: $\bar x \pm ME = 6.8 \pm $ \blank{2cm} $= ($ \blank{2cm} $,$ \blank{2cm} $)$ hours

\textbf{Step 4 --- Conclude:} \writelines{2}

\end{tcolorbox}

\begin{tcolorbox}[colback=white, colframe=black!40, title=\textbf{Tier A --- Approaching Proficiency},
    fonttitle=\bfseries, arc=2mm, left=3mm, right=3mm, top=2mm, bottom=2mm]

\textbf{Commute Times.} A city planner wants to estimate the mean commute time for
workers in a downtown district. A random sample of 36 workers is taken; summary
statistics: $\bar x = 28.4$ min, $s = 9.6$ min.

Construct a 95\% confidence interval for the true mean commute time. Use all four
steps and verify conditions explicitly.

\vspace{0.1in}
\textbf{Step 1:} \writeline

\textbf{Step 2 --- Conditions:} \writelines{4}

\textbf{Step 3:}

$df = $ \blank{2cm} \quad $t^* = $ \blank{2.5cm} \quad $SE = $ \blank{2.5cm}

Interval: \blank{10cm} $= ($ \blank{2cm} $,$ \blank{2cm} $)$ min

\textbf{Step 4:} \writelines{2}

\vspace{0.05in}
\textbf{Bonus:} How wide must the interval be? How could the planner halve the margin
of error? \writelines{2}

\end{tcolorbox}

\begin{tcolorbox}[colback=white, colframe=black!40, title=\textbf{Tier E --- Extension},
    fonttitle=\bfseries, arc=2mm, left=3mm, right=3mm, top=2mm, bottom=2mm]

\textbf{Matched-Pairs: Step-Count Intervention.} A fitness researcher measured daily
step counts for 12 participants before and after a 4-week intervention. The differences
(after $-$ before, in thousands of steps) are:

\medskip
\begin{center}
\begin{tabular}{cccccccccccc}
1.2 & 0.8 & $-0.3$ & 2.1 & 1.5 & 0.6 & 1.8 & $-0.1$ & 2.4 & 0.9 & 1.3 & 1.7
\end{tabular}
\end{center}
\medskip

Compute $\bar d$ and $s_d$, verify conditions, and construct a 95\% CI for the mean
change in daily step count. Then interpret the interval and assess whether the data
provide convincing evidence that the intervention increased step counts on average.

$\bar d = $ \blank{3cm} \quad $s_d = $ \blank{3cm} \quad $n = $ \blank{2cm}

\vspace{0.1in}
\textbf{Step 2 --- Conditions:} \writelines{3}

\textbf{Step 3:} $df = $ \blank{2cm} \quad $t^* = $ \blank{2.5cm}

Interval: $\bar d \pm t^* \cdot s_d/\sqrt{n} = $ \blank{8cm} $= ($ \blank{2cm} $,$ \blank{2cm} $)$ thousands of steps

\textbf{Step 4 and claim assessment:} \writelines{3}

\end{tcolorbox}

\end{document}
